\section{E1: the arms}
\label{sec:results}

\begin{table}[h]
\centering\small
\caption{CIFAR-10 test accuracy by arm. Mean $\pm$ standard deviation of the best test
accuracy over seeds. \code{clauses} is the total clause budget of the arm (both layers where
there are two), \code{TA} the automaton count in millions, and \code{train} the summed
per-epoch training time across all trainable stages on one RTX~3090.}
\label{tab:arms}
\pgfplotstabletypeset[
  col sep=comma, string type,
  every head row/.style={before row=\ttoprule, after row=\ttmidrule},
  every last row/.style={after row=\ttbotrule},
  columns={label,clauses,automata,acc_mean,acc_std,wall_min},
  columns/label/.style={column name=\thead{arm}, column type={l}},
  columns/clauses/.style={column name=\thead{clauses}, column type={r}},
  columns/automata/.style={column name=\thead{TA (M)}, column type={r}},
  columns/acc_mean/.style={column name=\thead{test acc (\%)}, column type={r}},
  columns/acc_std/.style={column name=\thead{$\pm$}, column type={r}},
  columns/wall_min/.style={column name=\thead{train (min)}, column type={r}},
]{data/arms.csv}
\end{table}

\begin{figure}[h]
\centering
\begin{tikzpicture}
\begin{axis}[
  width=0.92\textwidth, height=6.6cm,
  xlabel={epoch}, ylabel={test accuracy (\%)},
  grid=major, grid style={ttline!40},
  legend pos=south east, legend style={font=\sffamily\scriptsize, draw=ttline},
  label style={font=\sffamily\small}, tick label style={font=\sffamily\footnotesize},
  every axis plot/.append style={very thick},
]
\addplot[ttslate] table[col sep=comma, x=epoch, y=test_acc] {data/curve_single-l1.csv};
\addplot[ttslatel] table[col sep=comma, x=epoch, y=test_acc] {data/curve_single-matched.csv};
\addplot[ttmoss] table[col sep=comma, x=epoch, y=test_acc] {data/curve_single-rf.csv};
\addplot[ttamber] table[col sep=comma, x=epoch, y=test_acc] {data/curve_flat-stack.csv};
\addplot[ttorange] table[col sep=comma, x=epoch, y=test_acc] {data/curve_mctm.csv};
\addplot[ttorange, dashed] table[col sep=comma, x=epoch, y=test_acc] {data/curve_mctm-random.csv};
\legend{layer 1 alone, single matched, single 9x9, flat stack, MCTM, MCTM random L1}
\end{axis}
\end{tikzpicture}
\caption{Test accuracy per epoch (seed 0). For the two-layer arms the $x$ axis counts
layer-2 epochs; layer~1 is trained and frozen beforehand.}
\label{fig:curves}
\end{figure}

\subsection{What the second layer actually learned}
\label{sec:composition}

The density argument predicts a specific degeneracy: over a sparse input, a positive literal
(``channel $k$ fires here'') is rarely satisfied while its negation is almost always
satisfied, and Type~I feedback rewards including literals that are $1$ in the drawn patch.
A layer that has collapsed should therefore assert mostly \emph{absences}. Table~\ref{tab:comp}
measures this directly: \code{neg} is the median number of negated literals per clause and
\code{neg frac} the mean fraction of each clause's literals that are negations.

\begin{table}[h]
\centering\small
\caption{Clause composition of each arm's final (classifying) layer. \code{size} is the
median clause size, \code{pos}/\code{neg} the median count of plain and negated literals,
\code{neg frac} the mean negation fraction, and \code{empty} the fraction of clauses with no
included literals.}
\label{tab:comp}
\pgfplotstabletypeset[
  col sep=comma, string type,
  every head row/.style={before row=\ttoprule, after row=\ttmidrule},
  every last row/.style={after row=\ttbotrule},
  columns={label,size,pos,neg,neg_frac,empty},
  columns/label/.style={column name=\thead{arm}, column type={l}},
  columns/size/.style={column name=\thead{size}, column type={r}},
  columns/pos/.style={column name=\thead{pos}, column type={r}},
  columns/neg/.style={column name=\thead{neg}, column type={r}},
  columns/neg_frac/.style={column name=\thead{neg frac}, column type={r}},
  columns/empty/.style={column name=\thead{empty}, column type={r}},
]{data/composition.csv}
\end{table}

\subsection{Readout}

\lead{The stack is far worse than the layer it is built on.} The \mctm reaches
\accMctm{}\%, against \accSingleLOne{}\% for its own frozen layer~1 used directly as a
classifier. Adding a second layer does not fail to help; it removes more than half the
accuracy that was already there. The single-layer baselines at matched budget --- widening
to the same clause count (\accSingleMatched{}\%) or widening the patch to the stack's
$9\times9$ receptive field (\accSingleRf{}\%) --- are both far ahead.

\lead{Random first-layer clauses beat trained ones by a wide margin.} \code{mctm-random}
reaches \accMctmRandom{}\%, a gain of over twenty points on the identical architecture with a
\emph{trained} layer~1, and it is the only two-layer arm that is competitive with the
single-layer baselines at all. The difference is density and nothing else: random
three-literal clauses fire on \fireMctmRandom{}\% of positions, trained ones on
\fireMctm{}\%. Table~\ref{tab:comp} shows the consequence downstream --- the trained-layer-1
stack builds clauses that are \negMctm{}\% negations, the random-layer-1 stack
\negMctmRandom{}\%.

\lead{The ablations point the same way.} Removing position encoding from layer~2
(\accMctmNopos{}\%) and removing the OR-pool (\accMctmNopool{}\%) both leave the stack at or
near the $10\%$ chance floor, and the denser layer~1 of \code{mctm-dense}
(\accMctmDense{}\%) does not rescue it either. A $4\times4$ pool (\accMctmPoolFour{}\%) buys
density at the cost of a $15\times15$ receptive field, which is no longer comparable to the
$9\times9$ baseline.

\begin{ttcallout}[Read this with the scale in mind]
These are small Tsetlin machines --- hundreds of clauses where the CIFAR-10 literature uses
thousands --- so none of these absolute numbers is a CIFAR-10 result for Tsetlin machines.
What the table supports is the \emph{ordering} between arms at a matched clause budget, and
those gaps are an order of magnitude larger than the spread across seeds.
\end{ttcallout}

